\label{sec:data-sources}

All data sources for zone co-membership are free, publicly available from federal agencies,
and updated at least annually.
Table~\ref{tab:data-sources} summarizes the six zone types with their sources, formats,
update frequencies, and coverage.

\begin{table}[ht]
\centering
\caption{Zone co-membership data sources by priority. Coverage refers to the
  fraction of US states for which usable data exists. TIGER/Line and EIA
  Form 861 are the primary sources; all are free and public.}
\label{tab:data-sources}
\small
\begin{tabularx}{\textwidth}{@{}p{2.4cm}Xp{1.8cm}p{1.2cm}p{1.2cm}@{}}
\toprule
\textbf{Zone Type} & \textbf{Source / URL Pattern} & \textbf{Format} & \textbf{Update} & \textbf{Coverage} \\
\midrule
School district
  & TIGER/Line SCHOOLDISTRICT shapefile, census.gov/geo/
  & Shapefile
  & Annual
  & 100\% \\
County subdivision
  & TIGER/Line COUSUB shapefile, census.gov/geo/
  & Shapefile
  & Annual
  & 100\% \\
Electric utility territory
  & EIA Form 861 + HIFLD, eia.gov/electricity/data/eia861/
  & ZIP+Shapefile
  & Annual
  & 100\% \\
Fire district
  & TIGER/Line Special District (code 21), census.gov/geo/
  & Shapefile
  & Annual
  & $\sim$60\% \\
Water/sewer district
  & TIGER/Line Special District (codes 22--29), census.gov/geo/
  & Shapefile
  & Annual
  & $\sim$40\% \\
Police precinct
  & City/county GIS portals (no national source)
  & Varies
  & Varies
  & $\sim$30\% \\
\bottomrule
\end{tabularx}
\end{table}

\subsection{TIGER/Line School Districts}

The Census Bureau's TIGER/Line program publishes annual school district shapefiles
in three layers~\citep{census2020tiger}:
\begin{itemize}
  \item \textbf{Unified school districts} (UNSD): districts that provide both
    elementary and secondary education.
  \item \textbf{Elementary school districts} (ELSD): districts providing only
    elementary education.
  \item \textbf{Secondary school districts} (SCSD): districts providing only
    secondary education.
\end{itemize}
For census tracts in states where elementary and secondary functions are split
between different districts (as in parts of California, New Jersey, and New York),
co-membership is computed separately for each layer and averaged.
In states with unified districts (the majority), only the UNSD layer is needed.

The TIGER/Line school district shapefiles are published at
\url{https://www.census.gov/geographies/mapping-files/time-series/geo/tiger-line-file.html}
in both ESRI Shapefile and GeoJSON formats, with vintage years matching each decennial
census.
School district boundaries are stable between census years but change following
redistricting decisions by local school boards, which may trigger updates in the
annual TIGER/Line release.

\subsection{TIGER/Line County Subdivisions}

County subdivisions (COUSUB) are the Census Bureau's primary administrative unit
below the county level~\citep{census2020tiger}.
In New England states (Connecticut, Maine, Massachusetts, New Hampshire, Rhode Island,
Vermont), county subdivisions correspond to towns, which are the primary property
tax-levying unit.
In the Midwest (Michigan, Minnesota, Wisconsin, Indiana, and parts of Ohio, Kansas,
and the Dakotas), county subdivisions correspond to townships.
In the South and West, county subdivisions are less administratively meaningful (in
many southern counties they are simply ``census county divisions'' with no governmental
function), but they still provide a useful geographic unit for community co-membership
purposes.

COUSUB shapefiles are available from the same TIGER/Line distribution as school
district shapefiles, with 100\% national coverage.

\subsection{EIA Form 861 Electric Utility Service Territories}

The Energy Information Administration's Form 861 Annual Electric Power Industry Report
collects data from all electric utilities operating in the United States~\citep{eia861}.
The associated service territory shapefiles are available as a free annual ZIP
download from \url{https://www.eia.gov/electricity/data/eia861/}.

The EIA data covers all investor-owned utilities, municipal utilities, rural electric
cooperatives, and power marketing authorities.
Coverage is 100\% of US territory: every square meter of the contiguous United States
is served by exactly one retail electric utility (with the exception of small tribal
lands in a few states).

The HIFLD (Homeland Infrastructure Foundation-Level Data) Electric Retail Service
Territories dataset provides an alternative aggregation of the same EIA data at
\url{https://hifld-geoplatform.opendata.arcgis.com/datasets/electric-retail-service-territories}.
Both sources are updated annually and are equivalent for zone co-membership purposes.

\subsection{TIGER/Line Special Districts --- Fire and Water}

TIGER/Line Special Districts (the UNSD shapefile with type codes) include fire
protection districts (code 21) and various water and wastewater district types
(codes 22--29, covering water authorities, combined water/sewer, irrigation, and
drainage districts)~\citep{census2020tiger}.
Coverage for fire district data is approximately 60\% of states; water and sewer
district coverage is approximately 40\%.

The 60\% fire district coverage is not uniform: states in the Southeast (Georgia,
Alabama, Mississippi) and parts of the Midwest often rely on county or municipal
fire departments rather than special districts, and therefore have no TIGER/Line
fire district data.
Western states (California, Colorado, Oregon, Nevada, Washington) have extensive
fire district coverage.
The M.6 implementation checks for fire district data availability at the state level
before including it in $Z$.

\subsection{Police Precincts (Fallback to County)}

Municipal police precincts are available from city and county GIS portals for
approximately 30\% of the US population by coverage area.
There is no national data source for police precincts; the TIGER/Line program does
not maintain a precinct layer.

The M.6 implementation uses a two-level fallback:
\begin{enumerate}
  \item If a city GIS portal provides a precinct shapefile for the urbanized area
    containing the tract, use it.
  \item Otherwise, use county co-membership as a proxy for law enforcement co-membership
    (in rural areas, the county sheriff provides primary policing).
\end{enumerate}
Because county co-membership is already captured by the county subdivision dimension
(in states where county = county subdivision) or by county FIPS code directly,
this fallback does not introduce any new data source; it reuses existing county
boundary data.

\subsection{Spatial Join Implementation}

Zone assignment is performed by a centroid-in-polygon spatial join.
For each census tract:
\begin{enumerate}
  \item Compute the tract centroid (latitude/longitude of the population-weighted
    centroid from the TIGER/Line TRACT shapefile's INTPTLAT/INTPTLON fields).
  \item For each zone type $z$ in $Z$, find the polygon in the zone shapefile that
    contains the centroid and record its identifier.
  \item Store the zone assignment as a vector of identifiers for the tract.
\end{enumerate}
The spatial join is implemented using the Rust \texttt{geo} crate with a spatial
index (R-tree) for efficiency.
For a typical state with 1,000--5,000 tracts and six zone types, the spatial join
completes in under one second.

\subsection{Handling Zone Boundary Tracts}

Census tracts are designed by the Census Bureau to align with administrative boundaries
wherever possible~\citep{census2020tiger}.
In practice, however, some tracts straddle zone boundaries---particularly when zone
boundaries change between decennial censuses and tract boundaries have not been
updated to match.
For such tracts, the centroid assignment may place the tract in a zone that covers
only a minority of its population.

The implementation does not attempt to split tracts; it assigns each tract to exactly
one zone per zone type, determined by centroid.
Future work could improve accuracy by using areal interpolation (assigning the zone
that covers the largest fraction of tract area or population) or by aggregating to
block groups where tract-level ambiguity is detected.
For the current implementation, centroid assignment provides a conservative
approximation that errs on the side of underestimating co-membership (some tract
pairs will be assigned to different zones when their populations are mostly in the
same zone), which makes the zone co-membership signal conservative rather than
aggressive.
